\section{\DensityFB: normalized successor learning}
\label{sec:method}

\subsection{What is being predicted}

Write \( (s,a,s^+) \) for a transition from the offline dataset, and define
\begin{equation}
x=\phi(s), \qquad x^+=\phi(s^+)
\end{equation}
as the state variables whose future occupancy we want to predict.

The forward and backward representations see different inputs:
\begin{equation}
\label{eq:fb-args}
F(s,a,g)\in\mathbb R^d, \qquad B(x)\in\mathbb R^d.
\end{equation}
$F$ receives the full current state and action, while $B$ represents a candidate future configuration.

In our manipulation experiments, $\phi$ extracts the object coordinates.
We therefore condition on the full state but predict future occupancy over the goal-relevant object configuration.
This is the pushforward of the full-state successor measure and is sufficient for any reward defined only on those coordinates.
Setting $\phi$ to the identity instead gives $x=s$ and asks $B$ to represent the full state; we are testing this separately.

\subsection{From a successor measure to a probability distribution}

For a goal-conditioned policy $\pi_g(a\mid s)$, define the successor measure over $x$ as
\begin{equation}
    M^{\pi_g}(s,a,X) = \E_{\pi_g} \left[ \sum_{t\ge0} \gamma^t \ind{x_{t+1}\in X} \;\middle|\; s_0=s,a_0=a \right].
\end{equation}
Its total mass is
\begin{equation}
    M^{\pi_g}(s,a,\mathcal X) = \frac{1}{1-\gamma},
\end{equation}
independent of $(s,a,g)$.

We therefore normalize it:
\begin{equation}
\label{eq:normalized-sm}
    \nu^{\pi_g}(\cdot\mid s,a) = (1-\gamma) M^{\pi_g}(s,a,\cdot).
\end{equation}
Now
\begin{equation}
\nu^{\pi_g}(\mathcal X\mid s,a)=1.
\end{equation}

This normalization only rescales the usual successor-measure Bellman equation.

\begin{align}
    M^{\pi_g}(\cdot\mid s,a)    &= P_\phi(\cdot\mid s,a) + \gamma\,\E_{s^+,a^+} \left[ M^{\pi_g}(\cdot\mid s^+,a^+) \right], \\
    \nu^{\pi_g}(\cdot\mid s,a)  &= (1-\gamma)P_\phi(\cdot\mid s,a) + \gamma\,\E_{s^+,a^+} \left[ \nu^{\pi_g}(\cdot\mid s^+,a^+) \right],
\label{eq:normalized-bellman}
\end{align}
where the second line follows by multiplying the first by $1-\gamma$ and using $\nu=(1-\gamma)M$.
Thus the factor $1-\gamma$ appears only in the immediate term because we normalize the measure to unit mass; the recursive coefficient remains $\gamma$.
Consequently, for any $\nu_1,\nu_2$, the immediate term cancels and
\begin{equation}
    \|\mathcal T\nu_1-\mathcal T\nu_2\|_{\mathrm{TV}} \le \gamma\|\nu_1-\nu_2\|_{\mathrm{TV}},
\end{equation}
so the normalized Bellman operator has the same $\gamma$-contraction as the original successor-measure operator.

\subsection{Normalize the FB scores}

For a minibatch $\{(s_i,a_i,s_i^+)\}_{i=1}^n,$ we construct a query set from its observed successors (the cube positions)
\begin{equation}
U = \{u_j\}_{j=1}^n, \qquad u_j=\phi(s_j^+).
\end{equation}

For transition $i$ and candidate future configuration $u_j$, we compute
\begin{equation}
\ell_{ij} = \frac{ F(s_i,a_i,g_i)^\top B(u_j) }{\tau}.
\end{equation}
Standard FB interprets the bilinear quantity directly as a density.
\DensityFB{} instead treats it as a logit and normalizes across candidate successors:
\begin{equation}
\label{eq:density-softmax}
p_{ij} = \frac{ \exp(\ell_{ij}) }{ \sum_{k=1}^n \exp(\ell_{ik})}.
\end{equation}
Hence
\begin{equation}
p_{ij}\ge0, \qquad \sum_j p_{ij}=1.
\end{equation}

The vector $p_i$ represents the normalized future occupancy
$\nu^{\pi_{g_i}}(\cdot\mid s_i,a_i)$ on the sampled query set $U$.
The query set is redrawn on every update, so its size is a Monte-Carlo sample size.

\subsection{The hard Bellman target}

Equation~\ref{eq:normalized-bellman} gives the exact measure-level recursion.
For a sampled transition, the observed $x_i^+$ provides a sample from the immediate distribution $P_\phi(\cdot\mid s_i,a_i)$.
Because every batch successor is included in $U$, row $i$'s observed successor occupies slot $i$.

Define the one-hot vector
\begin{equation}
    e_{ij} = \ind{j=i}.
\end{equation}

Next, sample
\begin{equation}
a_i^+
\sim
\pi_{\bar\omega}(\cdot\mid s_i^+,g_i)
\end{equation}
and let $\bar p_i^+$ be the target-network prediction from $(s_i^+,a_i^+,g_i)$ evaluated over the same query set $U$.

The sampled Bellman target is then
\begin{equation}
\label{eq:hard-target}
    p_i^{\mathrm{tgt}} = \sg\left[ (1-\gamma)e_i + \gamma\bar p_i^+ \right].
\end{equation}
Both $e_i$ and $\bar p_i^+$ have unit mass, so
\begin{equation}
    \sum_j p_{ij}^{\mathrm{tgt}} = (1-\gamma)+\gamma = 1.
\end{equation}

We fit the predicted occupancy to this target with cross-entropy:
\begin{equation}
\label{eq:density-loss}
\mathcal L_{\mathrm{DensityFB}} = - \frac1n \sum_{i=1}^n \sum_{j=1}^n p_{ij}^{\mathrm{tgt}} \log p_{ij}.
\end{equation}

The update has a simple interpretation.
The observed successor receives $1-\gamma$ of the normalized occupancy, while the remaining $\gamma$ is inherited from the predicted occupancy after the next state.
Equations~\ref{eq:density-softmax}--\ref{eq:density-loss} define hard \DensityFB{}.

\subsection{What changed relative to FB?}

Hard \DensityFB{} keeps the forward--backward factorization and the policy-conditioned Bellman bootstrap.
The main change is how successor occupancy is represented.
Standard FB interprets $F(s,a,z)^\top B(x)$ directly as an unconstrained successor density.
\DensityFB{} instead treats the same bilinear quantity as a score and normalizes the scores across sampled successor queries:
\begin{equation}
    p_{ij} = \operatorname{softmax}_j \left( F(s_i,a_i,g_i)^\top B(u_j)/\tau \right).
\end{equation}
This makes the represented object the normalized successor measure $\nu=(1-\gamma)M$, whose sampled Bellman target is 
\begin{equation}
p_i^{\mathrm{tgt}} = (1-\gamma)e_i+\gamma\bar p_i^+.
\end{equation}

Two estimator choices accompany this change.
First, we fit the distributional target with cross-entropy rather than FB's squared/linear density estimator.
Second, the query set is built from batch successors, guaranteeing that the observed successor appears in the support of the immediate one-hot target.

We also remove FB's orthonormalization objective.
The softmax does not make $B$ orthogonal; rather, $F^\top B$ is no longer a calibrated density.
Only the resulting probabilities have predictive meaning, and downstream reward evaluation is performed directly under this distribution.
We therefore do not impose an orthonormal constraint on $B$.

\subsection{Why normalization addresses stationary play}

Normalization does not delete dwell occupancy.
If the cube remains stationary for many transitions, the learned distribution can still assign high probability to that configuration.
What it cannot do is increase its total mass:
\begin{equation}
    \sum_j p_{ij}=1.
\end{equation}
Repeated stationary transitions can therefore change where the predicted occupancy is allocated, but they cannot increase critic scale simply because the same successor appears many times.
This preserves the fact that the cube tends to remain where it was left while removing scale as a degree of freedom.

\subsection{Soft immediate targets}

Hard \DensityFB{} places the immediate $1-\gamma$ mass exactly on the observed successor.
We also evaluate a smoothed version.

Define normalized kernel weights
\begin{equation}
k_{ij}^{\epsilon} = \frac{ \exp\!\left( -\frac12 \|(u_j-x_i^+)/\epsilon\|^2 \right)}{ \sum_m \exp\!\left( -\frac12 \|(u_m-x_i^+)/\epsilon\|^2 \right)
}.
\end{equation}
Soft \DensityFB{} changes only the immediate term:
\begin{equation}
\label{eq:soft-target}
    p_{i,\mathrm{soft}}^{\mathrm{tgt}} = \sg\left[ (1-\gamma)k_i^\epsilon + \gamma\bar p_i^+\right].
\end{equation}

Everything else is unchanged.
The hard variant learns the normalized successor measure itself, whereas the soft variant smooths each one-step successor over an $\epsilon$-neighborhood before propagating occupancy.

\subsection{Role of the goal}

Goal conditioning is not required by the \DensityFB{} estimator.
We use $g$ because hindsight relabeling supplies goals from the play data, allowing us to learn a family of goal-conditioned policies $\{\pi_g\}$ and their corresponding successor measures.
In this sense, $g$ plays the same indexing role that the latent task vector $z$ plays in standard FB, and training still requires no reward labels or task labels.

\subsection{Control}

The policy-extraction recipe is unchanged from the corresponding baseline.
The only difference is the critic supplied to the actor.
A standard critic returns a scalar $Q(s,a,g)$ directly; \DensityFB{} returns a distribution over candidate future configurations, so we first read a scalar
goal value from that distribution.

For either hard or soft \DensityFB{}, define the goal readout
\begin{equation} 
    h_g^\epsilon(u_j) \propto \exp\!\left( -\frac12 \|(u_j-g)/\epsilon\|^2 \right), \qquad \sum_j h_g^\epsilon(u_j)=1.
\end{equation}
The scalar critic supplied to the actor is
\begin{equation}
\label{eq:goal-readout}
    q_g(s,a) = \sum_j p_j(s,a,g)\, h_g^\epsilon(u_j).
\end{equation}
The same goal readout is used for both training variants.
Hard versus soft refers only to the immediate term of the successor-measure
Bellman target, not to the actor or its goal weighting.

The occupancy model and the goal used to train the actor are separate objects.
For a goal $g$, the actor scores actions by the predicted occupancy near the goal, using the same normalized kernel weights as the soft immediate target, evaluated at the goal:
\begin{equation}
\label{eq:goal-readout}
q_g(s,a) = \sum_j p_j(s,a,g)\, k_j^\epsilon(g),
\end{equation}
with $k^\epsilon$ from Equation~\ref{eq:soft-target}'s kernel and the per-block tolerances for multi-object goals.
The kernel weighting is used by both training variants; the hard/soft choice affects only the immediate term of the measure loss, never the actor's goal weighting.

For the flow actor, we use this inside the behavior-regularized objective
\begin{equation}
\label{eq:actor}
    \mathcal L_\pi = \lambda_{\mathrm{BC}} \mathcal L_{\mathrm{BC\text{-}flow}} - \E\left[ \frac{ q_g(s,\pi(s,g)) }{s_q} \right],
\end{equation}
where
\begin{equation}
    s_q = \E\left[ \left| q_g-\E[q_g] \right| \right], \qquad \lambda_{\mathrm{BC}}=3.
\end{equation}
The predicted occupancy contains a large action-independent component, whereas the actor only needs differences between actions.
We therefore center $q_g$ before computing its magnitude normalizer.
The behavior-cloning baseline trains the same actor with the $q_g$ term removed.
We separately evaluate Gaussian DDPG+BC as an alternative policy extractor.

\subsection{One training step}
Algorithm~\ref{alg:densityfb} summarizes one \DensityFB{} update.

\begin{algorithm}[t]
\caption{
\DensityFB{}, one gradient step.
{\color{fbdiff}Blue} marks changes to the successor-measure estimator relative to standard FB.
{\color{hardvar}Teal}/{\color{softvar}orange} marks the hard/soft immediate-target choice.
Goal conditioning, the choice of $\phi$, and policy extraction are experimental setup choices and are left black.
}
\label{alg:densityfb}
\begin{algorithmic}[1]

    \State Sample transitions $\{(s_i,a_i,s_i^+)\}_{i=1}^{n}$ and hindsight goals $\{g_i\}_{i=1}^{n}$
    \State Set $x_i^+=\phi(s_i^+)$
    \State {\color{fbdiff} Build the query support $U=\{u_j=x_j^+\}_{j=1}^{n}$} \Comment{{\color{fbdiff}FB samples queries $x\sim\rho$}}
    \State Compute bilinear scores $c_{ij}=F(s_i,a_i,g_i)^\top B(u_j)$
    \State {\color{fbdiff} Normalize $p_{ij} = \operatorname{softmax}_j(c_{ij}/\tau)$} \Comment{{\color{fbdiff}FB treats $F^\top B$ as a density}}
    \State Sample bootstrap actions $a_i^+ \sim \pi_{\bar\omega}(\cdot\mid s_i^+,g_i)$
    \State Evaluate target-network occupancies $\bar p_i^+$ over the same query support $U$
    \State {\color{hardvar} Hard: $\eta_i=e_i$, where $e_{ij}=\ind{j=i}$} \Comment{{\color{softvar}Soft: $\eta_i=k_i^\epsilon$}}
    \State {\color{fbdiff} Form $p_i^{\mathrm{tgt}} = \sg\!\left[ (1-\gamma)\eta_i + \gamma\bar p_i^+ \right]$} \Comment{{\color{fbdiff}normalized Bellman target}}
    \State {\color{fbdiff} Update $F,B$ with $-\frac1n \sum_i\sum_j p_{ij}^{\mathrm{tgt}} \log p_{ij}$} \Comment{{\color{fbdiff}FB: squared/linear loss $+$ ortho}} 
    \State Update the actor using $q_g$ from Equation~\ref{eq:goal-readout}, with $F,B$ frozen \Comment{same extractor as baseline}
    \State Soft-update the target networks

\end{algorithmic}
\end{algorithm}
